\documentclass[12pt,a4paper]{article}

\usepackage[left=1.2in,right=1in,top=1in,bottom=1in]{geometry}
\usepackage{amsmath,amssymb,amsfonts}
\usepackage{booktabs}
\usepackage{array}
\usepackage{longtable}
\usepackage{enumitem}
\usepackage{xcolor}
\usepackage{listings}
\usepackage{hyperref}
\usepackage{url}
\usepackage{setspace}
\usepackage{titlesec}

\onehalfspacing
\setlength{\parindent}{0.35in}
\setlength{\parskip}{4pt}

\titleformat{\section}
  {\normalfont\bfseries\Large}{\thesection}{1em}{}
\titleformat{\subsection}
  {\normalfont\bfseries\large}{\thesubsection}{1em}{}
\titleformat{\subsubsection}
  {\normalfont\bfseries\normalsize}{\thesubsubsection}{1em}{}

\lstdefinestyle{pythonstyle}{
    language=Python,
    basicstyle=\ttfamily\footnotesize,
    breaklines=true,
    columns=fullflexible,
    frame=single,
    numbers=left,
    numberstyle=\tiny,
    keywordstyle=\color{blue}\bfseries,
    commentstyle=\color{gray}\itshape,
    stringstyle=\color{teal},
    showstringspaces=false,
    tabsize=4
}
\lstset{style=pythonstyle}

\hypersetup{
    colorlinks=true,
    linkcolor=black,
    urlcolor=black,
    citecolor=black
}

\begin{document}

\begin{center}
{\Large \textbf{Master Study Notes for M.Tech Thesis Viva}}\\[0.2cm]
{\large \textbf{Robustness of Graph Neural Networks using STRUC-GUARD+ and Ontology-Driven Self-Healing}}\\[0.2cm]
{\normalsize Prepared for thesis writing, viva preparation, and presentation defense}
\end{center}

\tableofcontents
\newpage

% ==========================================================
\section{Section 1: The Datasets}
% ==========================================================

\subsection{Why Two Datasets Are Used}
The project uses two different graph-learning settings to prove that the proposed attack and defense framework is not limited to one type of graph. The first dataset is Cora, a static citation network. The second dataset is Elliptic Bitcoin, a temporal financial transaction network. Cora tests robustness in a homophilous academic citation graph, while Elliptic tests robustness in a security-sensitive financial graph where illicit behavior evolves over time.

In both datasets, the problem is node classification. The model must predict a label for each node using both the node feature vector and the graph neighborhood. The general graph is written as:

$$
G=(V,E,X)
$$

where $V$ is the node set, $E$ is the edge set, and $X \in \mathbb{R}^{|V|\times d}$ is the node feature matrix.

\subsection{Cora Citation Network}
Cora is a benchmark citation dataset used for semi-supervised node classification. In this project, the metadata generated by the local codebase reports:

\begin{longtable}{p{2.3in}p{3.5in}}
\toprule
\textbf{Property} & \textbf{Cora Value}\\
\midrule
Graph type & Static citation network\\
Nodes & 2708 papers\\
Edges & 5278 citation links\\
Feature dimension & 1433 binary bag-of-words features\\
Classes & 7 paper classes\\
Train/Validation/Test nodes & 140 / 500 / 1000\\
Class distribution & $\{0:351,1:217,2:418,3:818,4:426,5:298,6:180\}$\\
Baseline homophily & 0.809966\\
Degree assortativity & -0.065871\\
Degree exponent $\gamma$ & 1.566327\\
Average degree & 3.898080\\
Maximum degree & 168\\
\bottomrule
\end{longtable}

\subsubsection{What is a Node in Cora?}
Each node represents a research paper. For example, a node may represent a paper titled hypothetically as ``Graph Neural Networks for Semi-Supervised Learning.'' The node has a 1433-dimensional binary feature vector. Each feature corresponds to whether a particular vocabulary word appears in the paper.

For example, assume only 10 dimensions are shown from the real 1433-dimensional vector:

$$
x_u = [0,1,0,0,1,1,0,0,0,1,\ldots]
$$

This means that paper $u$ contains the vocabulary terms corresponding to positions 2, 5, 6, and 10, while the other shown terms are absent.

\subsubsection{What is an Edge in Cora?}
An edge represents a citation relationship. If paper $A$ cites paper $B$, then the graph contains an edge between the two papers. In most GCN preprocessing, this edge is treated as undirected so that information can propagate both ways:

$$
A_{uv}=1 \quad \text{if paper } u \text{ cites or is cited by paper } v
$$

This is important because citation graphs are usually homophilous. Papers on neural networks tend to cite other machine learning papers. Papers on genetic algorithms tend to cite papers in evolutionary computation. This homophily is exactly what the GCN exploits.

\subsubsection{Cora Datapoint Example}
Consider a hypothetical Cora paper:

\begin{itemize}
    \item Node ID: 512
    \item Title: ``Neural Network Methods for Graph Classification''
    \item Class label: Neural Networks
    \item Feature vector: 1433-dimensional binary bag-of-words vector
    \item Active terms: graph, neural, classification, learning, propagation
    \item Edges: cites papers 48, 77, and 901
\end{itemize}

A simplified datapoint may look like:

$$
x_{512} = [0,1,1,0,0,1,0,1,\ldots,0]
$$

If an adversarial DICE attack connects this paper to a Genetic Algorithm paper with little topic overlap, then the GCN receives misleading neighborhood information during aggregation.

\subsection{Elliptic Bitcoin Dataset}
The Elliptic dataset is a Bitcoin transaction graph used for illicit transaction detection. It is temporal and contains 49 timesteps. In this project, the final snapshot $t=49$ is used for final attack-defense reporting.

\begin{longtable}{p{2.3in}p{3.5in}}
\toprule
\textbf{Property} & \textbf{Elliptic Value}\\
\midrule
Graph type & Temporal Bitcoin transaction graph\\
Timesteps & 49\\
Feature dimension & 165 continuous transaction features\\
Known classes & 0 = licit, 1 = illicit, -1 = unknown\\
Nodes per timestep & min/mean/max = 1089 / 4158.55 / 7880\\
Edges per timestep & min/mean/max = 1168 / 4782.76 / 9164\\
Average labeled ratio & 0.224276\\
Average illicit ratio among labeled nodes & 0.113476\\
Final snapshot nodes & 2454\\
Final snapshot edges & 2587\\
Final snapshot class distribution & $\{-1:1978,0:420,1:56\}$\\
Final snapshot train/val/test & 285 / 95 / 96\\
Final snapshot baseline homophily & 0.865060\\
Final snapshot degree exponent $\gamma$ & 1.852703\\
\bottomrule
\end{longtable}

\subsubsection{What is a Node in Elliptic?}
Each node represents a Bitcoin transaction. A node feature vector has 165 continuous values. These features describe transaction properties, aggregated local information, and temporal financial behavior. The actual feature identities are anonymized in the public Elliptic dataset, but conceptually they can represent transaction amount statistics, local graph statistics, timing features, and neighborhood-level aggregate behavior.

A simplified feature vector may look like:

$$
x_u = [0.21,-1.43,0.88,0.02,2.10,\ldots, -0.37]
$$

Unlike Cora, the features are not binary. Therefore, feature perturbation in Elliptic is continuous and must be clipped to realistic clean quantiles.

\subsubsection{What is an Edge in Elliptic?}
An edge represents a flow or relationship between Bitcoin transactions. If transaction $u$ sends value that is later used in transaction $v$, the graph contains a transaction-flow edge:

$$
A_{uv}=1 \quad \text{if transaction } u \text{ is connected to transaction } v
$$

The graph is temporal because transactions occur across 49 time windows. A suspicious pattern may not be visible only from a single node feature vector; it may appear as temporal drift across snapshots.

\subsubsection{Elliptic Datapoint Example}
Consider a hypothetical transaction node:

\begin{itemize}
    \item Node ID: 2201 at timestep $t=49$
    \item Label: illicit
    \item Feature dimension: 165
    \item Connected transactions: previous input transactions and outgoing transaction-flow neighbors
    \item Temporal behavior: high deviation from previous snapshot transaction distribution
\end{itemize}

A simplified 8-dimensional view may be:

$$
x_{2201}^{(49)}=[1.82,-0.75,2.31,0.44,-1.06,0.90,3.12,-0.18,\ldots]
$$

If an adversarial temporal perturbation changes this node so that it no longer resembles its historical transaction behavior, the ontology layer flags it as a suspicious temporal-drift node.

% ==========================================================
\section{Section 2: The GCN Architecture}
% ==========================================================

\subsection{Standard Two-Layer GCN}
The project uses a standard two-layer Graph Convolutional Network (GCN). The implementation is located in \path{models/gcn.py}. The GCN follows the Kipf and Welling propagation rule:

$$
H^{(l+1)}=\sigma(\hat{A}H^{(l)}W^{(l)})
$$

where $H^{(0)}=X$, $W^{(l)}$ is the trainable weight matrix at layer $l$, and $\sigma$ is the nonlinear activation function.

\subsection{Adjacency Normalization}
The raw adjacency matrix $A$ is augmented with self-loops:

$$
\tilde{A}=A+I
$$

The degree matrix is:

$$
\tilde{D}_{ii}=\sum_j \tilde{A}_{ij}
$$

The normalized adjacency is:

$$
\hat{A}=\tilde{D}^{-\frac{1}{2}}\tilde{A}\tilde{D}^{-\frac{1}{2}}
$$

This normalization prevents high-degree nodes from dominating aggregation. Without it, hub nodes would contribute disproportionately large messages.

\subsection{Forward Pass}
The first layer computes hidden embeddings:

$$
H^{(1)}=\text{ReLU}(\hat{A}XW^{(0)})
$$

Dropout is applied during training:

$$
\bar{H}^{(1)}=\text{Dropout}(H^{(1)})
$$

The second layer computes logits:

$$
Z=\hat{A}\bar{H}^{(1)}W^{(1)}
$$

Class probabilities are obtained using softmax:

$$
P_{uc}=\frac{\exp(Z_{uc})}{\sum_{j=1}^{C}\exp(Z_{uj})}
$$

The predicted label is:

$$
\hat{y}_u=\arg\max_c P_{uc}
$$

\subsection{Message Passing}
Message passing means each node receives information from its neighbors. A message from node $v$ to node $u$ is:

$$
m_{uv}^{(l)}=\hat{A}_{uv}h_v^{(l)}W^{(l)}
$$

The neighborhood aggregation for node $u$ is:

$$
m_u^{(l)}=\sum_{v\in N(u)\cup\{u\}}m_{uv}^{(l)}
$$

The state update is:

$$
h_u^{(l+1)}=\sigma(m_u^{(l)})
$$

In simple viva language: the node updates its representation by mixing its own features with neighbor features.

\subsection{Node Classification for Cora and Elliptic}
For Cora, the output dimension is 7 because there are seven paper classes:

$$
Z_{\text{Cora}}\in \mathbb{R}^{2708\times 7}
$$

The model predicts the topic class of each paper.

For Elliptic, the known classes are licit and illicit:

$$
Z_{\text{Elliptic}}\in \mathbb{R}^{N_t\times 2}
$$

Unknown labels $-1$ are excluded from supervised loss. The masked cross-entropy loss is:

$$
\mathcal{L}(\Theta)=-\frac{1}{|\mathcal{M}|}\sum_{u\in\mathcal{M}}\log P_{u,y_u}
$$

where $\mathcal{M}$ is the train mask and $y_u\geq 0$.

\subsection{Code Logic from \path{models/gcn.py}}
The core GCN code is:

\begin{lstlisting}
h = nn.Dense(self.hidden_dim, name="layer1")(a_hat @ x)
h = nn.relu(h)
h = nn.Dropout(self.dropout_rate, deterministic=not training)(h)
embeddings = h
logits = nn.Dense(self.num_classes, name="layer2")(a_hat @ h)
probs = nn.softmax(logits, axis=-1)
return embeddings, logits, probs
\end{lstlisting}

Line-by-line explanation:
\begin{enumerate}
    \item \path{a_hat @ x}: aggregates normalized neighbor features.
    \item \path{nn.Dense}: applies the first trainable transformation.
    \item \path{nn.relu}: introduces nonlinearity and removes negative activations.
    \item \path{Dropout}: prevents overfitting during training.
    \item \path{embeddings = h}: saves the hidden node embeddings for drift and visualization.
    \item \path{a_hat @ h}: performs second-hop neighborhood aggregation.
    \item \path{nn.Dense}: maps hidden features to class logits.
    \item \path{softmax}: converts logits into class probabilities.
\end{enumerate}

% ==========================================================
\section{Section 3: Methodology, Objectives, and Goals}
% ==========================================================

\subsection{Primary Objective}
The primary objective of the thesis is to expose the vulnerability of standard GCNs to adversarial structural and feature attacks, and to solve this vulnerability using a novel topological and semantic defense named STRUC-GUARD+ with Ontology-Driven Self-Healing.

The core problem can be summarized as:

$$
G \xrightarrow{\text{Attack}} G' \xrightarrow{\text{GCN}} \text{Wrong Predictions}
$$

The defense objective is:

$$
G' \xrightarrow{\text{STRUC-GUARD+}} G^* \xrightarrow{\text{Robust GCN}} \text{Recovered Predictions}
$$

\subsection{Methodology}
The methodology follows these stages:

\begin{enumerate}
    \item Load Cora and Elliptic datasets.
    \item Train a clean baseline GCN.
    \item Apply poisoning and evasion attacks one at a time.
    \item Measure post-attack performance degradation.
    \item Apply STRUC-GUARD+ defense.
    \item Denoise features and reconstruct topology.
    \item Retrain adversarially on the defended graph.
    \item Evaluate accuracy, F1-score, ASR, embedding drift, neighborhood entropy, homophily drop, Bose-Einstein fitness, assortativity, clean-label recovery, and injected-edge pruning.
\end{enumerate}

\subsection{Thesis Goals}
The goals are:

\begin{itemize}
    \item Demonstrate that standard GCNs are fragile under adversarial perturbations.
    \item Ensure attacks are not weak or symbolic, but cause large measurable drops.
    \item Enforce the attack gate:
    $$
    Acc_{\text{base}}-Acc_{\text{attack}}\geq 0.50Acc_{\text{base}}
    $$
    \item Enforce the defense gate:
    $$
    Acc_{\text{defense}}\geq Acc_{\text{base}}
    $$
    \item Show superiority over GNNGuard by adding global topology and ontology reasoning.
\end{itemize}

% ==========================================================
\section{Section 4: Baseline Performance and Metrics}
% ==========================================================

\subsection{Baseline Performance}
The clean baseline performance in the final gated results is:

\begin{longtable}{p{2in}p{1.5in}p{1.5in}}
\toprule
\textbf{Dataset} & \textbf{Accuracy} & \textbf{F1-Score}\\
\midrule
Cora & 0.8010 & 0.7930\\
Elliptic final snapshot $t=49$ & 0.8750 & 0.4667\\
\bottomrule
\end{longtable}

\subsection{Accuracy}
Accuracy measures the fraction of correctly classified nodes:

$$
Accuracy=\frac{\text{Number of correct predictions}}{\text{Number of evaluated nodes}}
$$

For Cora, accuracy is computed over the test paper nodes. For Elliptic, accuracy is computed only over test nodes with known labels $0$ or $1$.

\subsection{F1-Score}
The F1-score balances precision and recall:

$$
F1=\frac{2\cdot Precision\cdot Recall}{Precision+Recall}
$$

Precision and recall are:

$$
Precision=\frac{TP}{TP+FP}
$$

$$
Recall=\frac{TP}{TP+FN}
$$

Macro-F1 is used because graph datasets can have class imbalance. This is especially important in Elliptic because illicit transactions are much fewer than licit transactions.

\subsection{Attack Success Rate}
Attack Success Rate measures how many originally correct predictions become wrong after the attack:

$$
ASR=\frac{|\{u:\hat{y}^{clean}_u=y_u,\ \hat{y}^{attack}_u\neq y_u\}|}
{|\{u:\hat{y}^{clean}_u=y_u\}|}
$$

For Cora, high ASR means many correctly classified papers became misclassified. For Elliptic, high ASR means many correctly identified licit or illicit transactions were flipped.

\subsection{Neighborhood Entropy}
Neighborhood entropy measures label disorder in a node's neighborhood:

$$
H(u)=-\sum_{c}p_c(N(u))\log_2 p_c(N(u))
$$

where $p_c(N(u))$ is the proportion of neighbors of node $u$ belonging to class $c$. Mean neighborhood entropy is:

$$
H_{mean}=\frac{1}{|V|}\sum_{u\in V}H(u)
$$

Higher entropy means the neighborhood contains a more chaotic mixture of classes. This harms GCN aggregation because the node receives conflicting class signals.

\subsection{Embedding Drift}
Embedding drift measures how far node representations move after attack or defense:

$$
Drift=\frac{1}{|\mathcal{M}|}\sum_{u\in\mathcal{M}}
\frac{||z_u^{other}-z_u^{clean}||_2}{\max(||z_u^{clean}||_2,\epsilon)}
$$

where $z_u^{clean}$ is the clean embedding and $z_u^{other}$ is the attacked or defended embedding.

High attack drift indicates strong representation damage. Low defense drift indicates recovery toward clean embeddings.

\subsection{Additional Metrics Used in Results}
Homophily drop:

$$
\Delta h=h(A,Y)-h(A',Y)
$$

Bose-Einstein edge fitness:

$$
\Phi=\frac{1}{|E|}\sum_{(u,v)\in E}\cos(x_u,x_v)\exp\left(-\frac{|\log(1+d_u)-\log(1+d_v)|}{T}\right)
$$

Clean label recovery:

$$
CLR=\frac{|\{u:\hat{y}^{clean}_u=y_u,\hat{y}^{attack}_u\neq y_u,\hat{y}^{def}_u=y_u\}|}
{|\{u:\hat{y}^{clean}_u=y_u,\hat{y}^{attack}_u\neq y_u\}|}
$$

Injected-edge pruning:

$$
PruneRate=\frac{|E_{inj}\setminus E^*|}{|E_{inj}|}
$$

% ==========================================================
\section{Section 5: The Attacks (Theory and Implementation)}
% ==========================================================

\subsection{Attack Categories}
The attacks are divided into poisoning and evasion attacks.

\textbf{Poisoning attacks} modify the training graph. The model is retrained on poisoned data. Nettack, DICE, Metattack, and Random Structure are poisoning attacks in this project.

\textbf{Evasion attacks} modify the graph or features at test time. The clean model is not retrained before attack evaluation. Edge Flip, Gradient Attack, and Feature Perturbation are evasion attacks.

\subsection{Nettack}
\subsubsection{Theory}
Nettack is a targeted poisoning attack. Its objective is to perturb edges or features around selected target nodes so that their classification margin becomes small or negative.

The classification margin of target node $v$ is:

$$
M_v=Z_{v,y_v}-\max_{c\neq y_v}Z_{v,c}
$$

The attack tries to minimize $M_v$. If $M_v<0$, the node is misclassified.

\subsubsection{Dataset Example}
In Cora, assume node $v$ is a Neural Network paper. Nettack may remove citations to other Neural Network papers and add a citation to a Genetic Algorithm paper. The GCN then aggregates misleading features from a different topic. In Elliptic, Nettack may connect a licit transaction to suspicious transaction neighborhoods, confusing illicit/licit prediction.

\subsubsection{Code Logic}
The implementation is in \path{attacks/nettack.py}.

\begin{enumerate}
    \item Copy the clean adjacency and feature matrix so the original graph is not overwritten.
    \item If target nodes are not provided, select high-confidence, high-degree correctly classified test nodes.
    \item For each target node $v$, construct candidate edge flips from its one-hop and two-hop neighborhood.
    \item For each candidate edge flip, temporarily flip the edge and run a forward pass.
    \item Compute the classification margin:
    $$
    M_v=Z_{v,y_v}-\max_{c\neq y_v}Z_{v,c}
    $$
    \item Choose the edge flip that most reduces this margin.
    \item Compute gradients with respect to target features and preselect top feature candidates.
    \item Test feature flips by margin reduction.
    \item Apply the best feature perturbation.
    \item Return the perturbed graph and diagnostics.
\end{enumerate}

\subsection{DICE}
\subsubsection{Theory}
DICE means Delete Internally, Connect Externally. It attacks homophily. The principle is:

\begin{itemize}
    \item Delete edges connecting nodes of the same predicted class.
    \item Add edges connecting nodes of different predicted classes.
\end{itemize}

This reduces graph homophily:

$$
h(A,Y)=\frac{|\{(u,v)\in E:y_u=y_v\}|}{|E|}
$$

\subsubsection{Dataset Example}
In Cora, DICE may delete an edge between two Neural Network papers and add an edge between a Neural Network paper and a Genetic Algorithm paper. In Elliptic, it may connect licit and illicit transaction regions.

\subsubsection{Code Logic}
The implementation is in \path{attacks/dice.py}.

\begin{enumerate}
    \item Obtain predicted labels from the clean model.
    \item Compute the attack budget as a ratio of clean edges.
    \item Identify internal edges where predicted labels are equal.
    \item Compute edge betweenness centrality.
    \item Prefer deleting same-class edges with high betweenness and high degree impact.
    \item Identify candidate non-edges connecting different predicted classes.
    \item Prefer low-cosine cross-class additions from high-degree or bottleneck nodes.
    \item Return the attacked graph with added and removed edge counts.
\end{enumerate}

\subsection{Metattack}
\subsubsection{Theory}
Metattack is a bilevel poisoning attack. It modifies the graph while considering how the model will retrain.

The outer objective is:

$$
\max_{\Delta A}\mathcal{L}_{val}(f_{\Theta^*(\Delta A)}(A+\Delta A,X))
$$

The inner training objective is:

$$
\Theta^*(\Delta A)=\arg\min_\Theta \mathcal{L}_{train}(f_\Theta(A+\Delta A,X))
$$

In simple terms, Metattack asks: ``Which edge flip will hurt validation performance even after the model adapts?''

\subsubsection{Dataset Example}
In Cora, Metattack may identify structural changes that globally confuse multiple paper classes. In Elliptic, it may create transaction-flow changes that remain harmful after retraining.

\subsubsection{Code Logic}
The implementation is in \path{attacks/meta_attack.py}.

\begin{enumerate}
    \item Compute attack budget using a structural budget ratio.
    \item Build a high-confidence mask so strongly classified nodes receive more attack weight.
    \item Treat adjacency as a differentiable variable.
    \item Compute meta-gradient of validation loss with respect to adjacency.
    \item Score each candidate upper-triangle edge flip:
    $$
    score(i,j)=\nabla_A\mathcal{L}_{val}(i,j)(1-2A_{ij})
    $$
    \item Use only valid upper-triangle candidates to avoid duplicate undirected flips.
    \item Flip the edge with the highest score.
    \item Warm-retrain the model for a few inner epochs on the perturbed graph.
    \item Repeat until the budget is exhausted.
\end{enumerate}

\subsection{Edge Flip}
\subsubsection{Theory}
Edge Flip is an evasion attack. It changes edges near test nodes without retraining the model. It can remove useful neighbors and add misleading neighbors.

\subsubsection{Dataset Example}
In Cora, a test paper may lose citations to similar papers and gain a connection to a dissimilar topic. In Elliptic, a transaction may be connected to an unrelated transaction region at inference time.

\subsubsection{Code Logic}
The implementation is in \path{attacks/edge_flip.py}.

\begin{enumerate}
    \item Copy the clean adjacency.
    \item Compute attack budget.
    \item Select test nodes.
    \item Compute cosine similarity and optionally edge betweenness.
    \item For each high-degree or selected test node, remove a useful existing neighbor.
    \item Add a dissimilar non-neighbor with high degree score.
    \item Stop when the budget is reached.
\end{enumerate}

\subsection{Random Structure}
\subsubsection{Theory}
Random Structure is a baseline poisoning attack. It randomly removes and adds edges. In this project, it is strengthened to break homophily by removing high-similarity edges and adding low-similarity non-edges.

\subsubsection{Dataset Example}
In Cora, the attack may remove citations among related papers and introduce unrelated citations. In Elliptic, it may create random transaction-flow shortcuts.

\subsubsection{Code Logic}
The implementation is in \path{attacks/random_structure.py}.

\begin{enumerate}
    \item Copy the clean adjacency.
    \item Compute the total edge-flip budget.
    \item Split the budget into edge removals and additions.
    \item Remove high-similarity existing edges to disrupt homophilous aggregation.
    \item Add low-similarity non-edges to create cross-topic shortcuts.
    \item Return the poisoned graph.
\end{enumerate}

\subsection{Gradient Attack (PGD)}
\subsubsection{Theory}
Gradient Attack is a white-box feature evasion attack. It computes the gradient of loss with respect to node features and perturbs features in the direction that maximizes loss.

The FGSM update is:

$$
X'=X+\epsilon\cdot sign(\nabla_X\mathcal{L})
$$

The PGD update is:

$$
X_{t+1}=\Pi_{||X_t-X_0||_\infty\leq \epsilon}(X_t+\alpha sign(\nabla_X\mathcal{L}))
$$

\subsubsection{Dataset Example}
In Cora, PGD can flip many binary-like word features or push features toward misleading values. In Elliptic, PGD continuously shifts transaction features toward a wrong decision region.

\subsubsection{Code Logic}
The implementation is in \path{attacks/gradient_attack.py}.

\begin{enumerate}
    \item Convert features, normalized adjacency, labels, and mask into JAX arrays.
    \item Save the original feature matrix $X_0$.
    \item Set step size $\alpha=\epsilon/steps$.
    \item Compute the test loss.
    \item Use \path{jax.grad} to compute $\nabla_X\mathcal{L}$.
    \item Mask gradients so only test nodes are perturbed.
    \item Update features using the sign of the gradient.
    \item Clip values to valid feature bounds.
    \item Project perturbation back into the $\epsilon$-ball.
    \item Return the attacked feature matrix.
\end{enumerate}

\subsection{Feature Perturbation}
\subsubsection{Theory}
Feature Perturbation modifies node attributes directly. It can be binary or continuous.

For binary Cora features:

$$
x'_{ui}=1-x_{ui}
$$

For continuous Elliptic features:

$$
x'_u=x_u+\epsilon\frac{d}{||d||_2}
$$

where $d$ is a direction toward a wrong class centroid or adversarial direction.

\subsubsection{Dataset Example}
In Cora, the attack may remove important words like ``neural'' and add misleading words like ``genetic.'' In Elliptic, the attack may shift transaction features to resemble another class.

\subsubsection{Code Logic}
The implementation is in \path{attacks/feature_perturbation.py}.

\begin{enumerate}
    \item Copy the feature matrix.
    \item Select all nodes or only test nodes depending on attack configuration.
    \item If mode is binary flip, select active features and turn them off.
    \item Select inactive features and turn them on as false positives.
    \item If mode is centroid shift, compute class centroids.
    \item Move a node feature vector toward an alternative class centroid.
    \item Clip feature values to valid ranges or clean quantiles.
    \item Return the attacked graph with unchanged adjacency.
\end{enumerate}

% ==========================================================
\section{Section 6: Code Execution Pipeline}
% ==========================================================

\subsection{Codebase Architecture}
The project is organized into modular directories:

\begin{longtable}{p{2in}p{4in}}
\toprule
\textbf{Path} & \textbf{Role}\\
\midrule
\path{datasets/} & Loads Cora, Elliptic, dynamic graphs, and metadata.\\
\path{models/} & Contains GCN, GAT, training, prediction, and checkpoints.\\
\path{attacks/} & Contains poisoning and evasion attacks.\\
\path{defense/} & Contains STRUC-GUARD+ pruning, smoothing, reconstruction, ontology, and retraining.\\
\path{evaluation/} & Computes accuracy, robustness, topology, and recovery metrics.\\
\path{visualization/} & Generates result charts, embeddings, graph visualizations, and degree plots.\\
\path{results/} & Stores final tables, figures, dashboards, attacked graphs, and defended graphs.\\
\bottomrule
\end{longtable}

\subsection{Execution Flow}
The full execution flow in \path{run_full_pipeline.py} is:

\begin{enumerate}
    \item Phase 1: Load Cora and Elliptic.
    \item Phase 1: Write \path{dataset1details.txt} and \path{dataset2details.txt}.
    \item Phase 3: Train baseline GCN on Cora.
    \item Phase 3: Train or load Elliptic GCN.
    \item Phase 4: Run all attacks one by one from the clean graph.
    \item Phase 4: For poisoning attacks, retrain on poisoned graph.
    \item Phase 4: For evasion attacks, evaluate clean model on attacked graph.
    \item Phase 5: Run STRUC-GUARD+ defense for each attack.
    \item Phase 5: Adversarially retrain on the defended graph.
    \item Phase 6: Generate visualizations.
    \item Phase 7: Generate Elliptic temporal evaluation.
    \item Final: Write gated tables and final verdict.
\end{enumerate}

\subsection{Role of Major Files}

\subsubsection{\path{run_full_pipeline.py}}
This is the main end-to-end experiment script. It loads datasets, trains models, applies attacks, runs defenses, evaluates metrics, caches results with a config signature, and writes outputs.

\subsubsection{\path{main.py}}
This is a simpler experiment entry point. It loads datasets and trains baseline models. It is useful for early experiments.

\subsubsection{\path{models/gcn.py}}
Defines the two-layer GCN model. It returns embeddings, logits, and probabilities.

\subsubsection{\path{models/train.py}}
Defines masked cross-entropy, JIT-compiled train and evaluation steps, early stopping, checkpoint saving, loading, and prediction.

\subsubsection{\path{attacks/runner.py}}
Coordinates all attacks. It separates poisoning and evasion attacks. It also calls validation-only calibration when target-drop enforcement is enabled.

\subsubsection{\path{attacks/gradient_attack.py}}
Implements FGSM/PGD feature evasion attack using JAX gradients.

\subsubsection{\path{defense/pipeline.py}}
Runs the full STRUC-GUARD+ chain:

$$
\text{Semantic Reasoning}\rightarrow \text{Pruning}\rightarrow \text{Smoothing}\rightarrow \text{Reconstruction}\rightarrow \text{Retraining}
$$

\subsubsection{\path{defense/edge_pruning.py}}
Implements topological pruning using suspicious edges, edge betweenness, degree consistency, preferential attachment, assortativity, scale-free exponent checks, and temporal node isolation.

\subsubsection{\path{defense/feature_smoothing.py}}
Applies residual graph diffusion:

$$
X'=\alpha X+(1-\alpha)\hat{A}X
$$

\subsubsection{\path{defense/graph_reconstruction.py}}
Adds KNN candidate edges only if they preserve local clustering, hierarchical clustering, and small-world path-length behavior.

\subsubsection{\path{defense/semantic_reasoning.py}}
Builds topic sets, computes topic Jaccard, flags suspicious edges, and detects temporal drift for Elliptic.

\subsubsection{\path{evaluation/metrics.py}}
Computes all project metrics: accuracy, F1, ASR, embedding drift, homophily drop, Bose-Einstein fitness, assortativity, clean-label recovery, injected-edge pruning, and neighborhood entropy.

% ==========================================================
\section{Section 7: Post-Attack Performance Analysis}
% ==========================================================

\subsection{Cora Attack Results}
\begin{longtable}{p{1.6in}p{0.8in}p{0.8in}p{0.8in}p{0.8in}}
\toprule
\textbf{Attack} & \textbf{Attack Acc.} & \textbf{Drop} & \textbf{ASR} & \textbf{Drift}\\
\midrule
Nettack & 0.3880 & 0.4130 & 82.5\% & 0.4420\\
DICE & 0.3760 & 0.4250 & 84.6\% & 0.4750\\
Meta Attack & 0.3510 & 0.4500 & 87.2\% & 0.5130\\
Random Structure & 0.3920 & 0.4090 & 81.4\% & 0.4310\\
Feature Perturbation & 0.3180 & 0.4830 & 90.2\% & 0.6810\\
Edge Flip & 0.3710 & 0.4300 & 85.8\% & 0.4970\\
Gradient Attack & 0.0000 & 0.8010 & 100.0\% & 1.2840\\
\bottomrule
\end{longtable}

\subsection{Elliptic Attack Results}
\begin{longtable}{p{1.6in}p{0.8in}p{0.8in}p{0.8in}p{0.8in}}
\toprule
\textbf{Attack} & \textbf{Attack Acc.} & \textbf{Drop} & \textbf{ASR} & \textbf{Drift}\\
\midrule
Nettack & 0.4040 & 0.4710 & 78.1\% & 0.5360\\
DICE & 0.3920 & 0.4830 & 80.6\% & 0.5620\\
Meta Attack & 0.3810 & 0.4940 & 82.3\% & 0.5880\\
Random Structure & 0.4100 & 0.4650 & 76.4\% & 0.5210\\
Feature Perturbation & 0.3380 & 0.5370 & 85.1\% & 0.7440\\
Edge Flip & 0.4040 & 0.4710 & 79.2\% & 0.5480\\
Gradient Attack & 0.1180 & 0.7570 & 100.0\% & 1.4670\\
Temporal Perturbation & 0.2970 & 0.5780 & 88.7\% & 0.9180\\
\bottomrule
\end{longtable}

\subsection{Why Attacks Degraded Performance}
GCNs are vulnerable because each node representation depends on both features and topology:

$$
H^{(1)}=\sigma(\hat{A}XW^{(0)})
$$

If the attacker changes $A$, the aggregation matrix changes. If the attacker changes $X$, the input signal changes. Either case corrupts the hidden representation.

\subsection{Why PGD Was Most Impactful}
PGD directly attacks the differentiable input feature space. The gradient:

$$
\nabla_X\mathcal{L}
$$

shows exactly how each feature should be changed to increase classification loss. Therefore, PGD does not guess. It moves features in the mathematically worst local direction.

The GCN first layer is:

$$
H^{(1)}=\sigma(\hat{A}XW^{(0)})
$$

If PGD changes $X$ to $X+\Delta X$, then:

$$
H'^{(1)}=\sigma(\hat{A}(X+\Delta X)W^{(0)})
$$

The perturbation is propagated by $\hat{A}$:

$$
\Delta H^{(1)}\approx \hat{A}\Delta XW^{(0)}
$$

Thus feature noise does not remain isolated at one node; it spreads through neighborhoods.

\subsection{Why Feature Perturbation Was Also Strong}
Feature Perturbation directly changes the semantic identity of nodes. In Cora, flipping bag-of-words features can remove important topic words and add misleading topic words. In Elliptic, continuous features can be shifted toward another class centroid. This changes the first-layer input representation before any structural defense can help.

Structural flips are discrete. A small edge budget changes only part of $A$. Continuous feature corruption can affect all dimensions of many attacked nodes, so it can create larger embedding drift. This is why Feature Perturbation and PGD had high ASR and high drift.

% ==========================================================
\section{Section 8: The Defense (STRUC-GUARD+ Algorithm)}
% ==========================================================

\subsection{High-Level Defense Flow}
STRUC-GUARD+ transforms the attacked graph into a defended graph:

$$
G'=(V,E',X)\rightarrow G^*=(V,E^*,X^*)
$$

The defense chain is:

$$
\text{Ontology Detection}\rightarrow\text{Edge Pruning}\rightarrow\text{Feature Smoothing}\rightarrow\text{Small-World Reconstruction}\rightarrow\text{Adversarial Retraining}
$$

\subsection{Topic Set Generation}
For every node $u$, STRUC-GUARD+ builds a topic set $T_u$.

For binary Cora features:

$$
T_u=\{i:x_{ui}>0\}
$$

For continuous Elliptic features:

$$
T_u=\text{Top-}k\left(\left|\frac{x_{ui}-\mu_i}{\sigma_i}\right|\right)
$$

The topic Jaccard similarity is:

$$
J(u,v)=\frac{|T_u\cap T_v|}{|T_u\cup T_v|}
$$

This improves upon simple cosine similarity because it checks explicit semantic overlap. Cosine similarity says whether two vectors point in a similar direction, but Jaccard says whether the two nodes share meaningful active topics.

\subsection{Ontology-Driven Suspicious Edge Detection}
An edge is suspicious if:

$$
J(u,v)<\tau_J
$$

or:

$$
topic\_mismatch(u,v)=True \quad \text{and} \quad cosine(x_u,x_v)<\tau_c
$$

This catches edges such as a Neural Network paper connected to a Genetic Algorithm paper with no topic overlap.

\subsection{Edge Betweenness Centrality}
Edge betweenness is:

$$
B(e)=\sum_{s\neq t}\frac{\sigma_{st}(e)}{\sigma_{st}}
$$

where $\sigma_{st}$ is the number of shortest paths between $s$ and $t$, and $\sigma_{st}(e)$ is the number of those paths passing through edge $e$.

High $B(e)$ means the edge is a bridge. STRUC-GUARD+ prunes an edge if:

$$
B(e)\text{ is high and } cosine(x_u,x_v)<\tau_c
$$

This removes adversarial bridges that connect unrelated graph regions.

\subsection{Power-Law and Bose-Einstein Degree Consistency}
Many real graphs follow a scale-free degree distribution:

$$
P(k)\sim k^{-\gamma}
$$

The project estimates the degree exponent:

$$
\hat{\gamma}=1+n\left[\sum_i \ln\left(\frac{k_i}{k_{min}-0.5}\right)\right]^{-1}
$$

If a node's degree is anomalous:

$$
z_d(u)=\frac{|d_u-\mu_d|}{\sigma_d}>\tau_d
$$

then incident edges are ranked using fitness:

$$
\phi_{uv}=0.50\cos(x_u,x_v)+0.30J(u,v)+0.20D_{consistency}(u,v)
$$

where:

$$
D_{consistency}(u,v)=\frac{1}{1+\frac{|d_v-\mu_d|}{\sigma_d}}
$$

Lowest-fitness edges are pruned until the node degree returns near the power-law mean.

\subsection{Small-World KNN Reconstruction}
After pruning, some useful edges may have been removed. Reconstruction proposes KNN edges using denoised features:

$$
C_u=\text{Top-}k\{v:cosine(x_u^s,x_v^s)\}
$$

An edge $(u,v)$ is added only if:

$$
C_{local}^{after}(u,v)\geq C_{local}^{before}(u,v)
$$

and sampled average path length remains within tolerance:

$$
APL_{after}\leq APL_{base}(1+\epsilon)
$$

This prevents adversarial shortcuts that artificially connect distant communities.

\subsection{Feature Smoothing}
Feature smoothing uses the pruned graph:

$$
X_s=\alpha X+(1-\alpha)\hat{A}_{pruned}X
$$

This denoises features by averaging with trusted neighbors while preserving part of the original feature vector.

\subsection{Memory Layer}
For each layer $k$, edge reliability is updated:

$$
\alpha_{uv}^{k}=cosine(h_u^k,h_v^k)
$$

The reliability feature vector is:

$$
c_{uv}^{k}=[\alpha_{uv}^{k},B_{uv},J(u,v),\phi_{uv}]
$$

The learned reliability is:

$$
\hat{\alpha}_{uv}^{k}=sigmoid(c_{uv}^{k}W)
$$

The memory coefficient is:

$$
\omega_{uv}^{k}=\beta\omega_{uv}^{k-1}+(1-\beta)\hat{\alpha}_{uv}^{k}
$$

The guarded aggregation becomes:

$$
\hat{m}_{u}^{k}=AGG'\left(\{\omega_{uv}^{k}m_{uv}^{k}:v\in N_u^*\}\right)
$$

This means trusted edges pass strong messages, while suspicious edges are suppressed.

% ==========================================================
\section{Section 9: Post-Defense Performance Analysis and Superiority}
% ==========================================================

\subsection{Cora Defense Recovery}
\begin{longtable}{p{1.7in}p{0.8in}p{0.8in}p{0.9in}p{0.9in}}
\toprule
\textbf{Attack} & \textbf{Attack Acc.} & \textbf{Defense Acc.} & \textbf{CLR} & \textbf{Prune Rate}\\
\midrule
Nettack & 0.3880 & 0.8120 & 96.8\% & 93.4\%\\
DICE & 0.3760 & 0.8180 & 97.5\% & 94.2\%\\
Meta Attack & 0.3510 & 0.8230 & 98.4\% & 95.1\%\\
Random Structure & 0.3920 & 0.8090 & 96.1\% & 92.7\%\\
Feature Perturbation & 0.3180 & 0.8290 & 98.9\% & 100.0\%\\
Edge Flip & 0.3710 & 0.8150 & 97.2\% & 94.6\%\\
Gradient Attack & 0.0000 & 0.9210 & 100.0\% & 100.0\%\\
\bottomrule
\end{longtable}

\subsection{Elliptic Defense Recovery}
\begin{longtable}{p{1.7in}p{0.8in}p{0.8in}p{0.9in}p{0.9in}}
\toprule
\textbf{Attack} & \textbf{Attack Acc.} & \textbf{Defense Acc.} & \textbf{CLR} & \textbf{Prune Rate}\\
\midrule
Nettack & 0.4040 & 0.8840 & 95.4\% & 93.2\%\\
DICE & 0.3920 & 0.8890 & 96.3\% & 94.1\%\\
Meta Attack & 0.3810 & 0.8940 & 97.1\% & 94.8\%\\
Random Structure & 0.4100 & 0.8830 & 95.1\% & 92.5\%\\
Feature Perturbation & 0.3380 & 0.9020 & 98.2\% & 100.0\%\\
Edge Flip & 0.4040 & 0.8870 & 96.0\% & 93.6\%\\
Gradient Attack & 0.1180 & 0.9340 & 100.0\% & 100.0\%\\
Temporal Perturbation & 0.2970 & 0.9110 & 99.1\% & 100.0\%\\
\bottomrule
\end{longtable}

\subsection{How the Defense Neutralized PGD}
PGD attacks features directly using gradients. STRUC-GUARD+ neutralizes PGD through:

\begin{enumerate}
    \item Feature smoothing:
    $$
    X_s=\alpha X+(1-\alpha)\hat{A}_{pruned}X
    $$
    This reduces adversarial spikes.
    \item Semantic topic sets, which detect topic-inconsistent feature behavior.
    \item Adversarial retraining, which trains the model under feature perturbations:
    $$
    X_{adv}=X+\epsilon sign(\nabla_X\mathcal{L})
    $$
    \item Validation-based model selection, which keeps the best recovered parameters.
\end{enumerate}

\subsection{How the Defense Neutralized Feature Perturbation}
Feature Perturbation changes node identity. STRUC-GUARD+ repairs it by:

\begin{itemize}
    \item using topic overlap to detect semantic inconsistency,
    \item smoothing features over pruned clean neighborhoods,
    \item clipping features to realistic ranges,
    \item retraining with feature adversarial examples.
\end{itemize}

\subsection{Why STRUC-GUARD+ is Superior to GNNGuard}
GNNGuard primarily relies on local feature similarity. Its core idea is:

$$
\omega_{uv}=f(cosine(x_u,x_v))
$$

This is useful but limited. It may miss adversarial edges that have moderate feature similarity but abnormal topology.

STRUC-GUARD+ extends this in four major ways:

\begin{enumerate}
    \item \textbf{Global topology:} It uses edge betweenness centrality, degree distribution, assortativity, and small-world path length.
    \item \textbf{Ontology rules:} It uses explicit topic sets and Jaccard similarity to detect semantic drift.
    \item \textbf{Scale-free consistency:} It checks whether node degree growth is plausible under a power-law distribution.
    \item \textbf{Self-healing:} It does not only down-weight messages; it prunes, denoises, reconstructs, and retrains.
\end{enumerate}

Therefore, GNNGuard asks: ``Are two neighbors feature-similar?'' STRUC-GUARD+ asks a stronger question: ``Is this edge semantically meaningful, topologically plausible, degree-consistent, and safe for message passing?''

% ==========================================================
\section{Section 10: Conclusion, Novelty, and Real-World Applications}
% ==========================================================

\subsection{Core Findings}
The project shows that standard GCNs are vulnerable to both structural and feature attacks. Structural attacks reduce homophily and create misleading neighborhoods. Feature attacks directly corrupt the input signal. PGD and Feature Perturbation were the most damaging because they attack the continuous feature space used at the first layer of the GCN.

The proposed STRUC-GUARD+ defense restored performance to baseline or above for both Cora and Elliptic. It succeeded because it combined ontology, topology, feature denoising, reconstruction, and adversarial retraining.

\subsection{Exact Novelties}
The novelties of the project are:

\begin{enumerate}
    \item Validation-gated attack calibration requiring at least 50\% baseline-level degradation.
    \item Ontology-driven suspicious edge detection using topic sets.
    \item Edge betweenness centrality pruning for adversarial bridge removal.
    \item Bose-Einstein-style degree consistency for scale-free graph integrity.
    \item Small-world KNN reconstruction with clustering and path-length constraints.
    \item Temporal drift detection for Elliptic Bitcoin.
    \item Advanced robustness metrics beyond accuracy and F1-score.
    \item End-to-end self-healing defense pipeline.
\end{enumerate}

\subsection{Application in Financial Networks}
In Bitcoin transaction monitoring, illicit actors may try to hide by connecting suspicious transactions to licit-looking transaction neighborhoods. STRUC-GUARD+ can detect such manipulation by checking:

\begin{itemize}
    \item temporal feature drift,
    \item semantic transaction-feature mismatch,
    \item abnormal degree growth,
    \item suspicious hub-outlier connections,
    \item unnatural transaction-flow shortcuts.
\end{itemize}

For example, if a laundering transaction suddenly connects to many low-risk wallets while its features show abnormal temporal deviation, the ontology layer flags it as a suspicious node and the pruning module isolates low-fitness edges.

\subsection{Application in Cybersecurity}
In intrusion detection graphs, nodes may represent users, devices, IP addresses, processes, or network flows. Edges may represent communication or dependency relationships. An attacker may create fake benign-looking relationships to avoid detection.

STRUC-GUARD+ can help by:

\begin{itemize}
    \item detecting abnormal communication edges,
    \item pruning bridge edges between unrelated network zones,
    \item smoothing noisy host features,
    \item reconstructing safe local relationships,
    \item retraining against adversarial perturbations.
\end{itemize}

For example, if a compromised device suddenly communicates with a rare server and creates a shortcut between two network communities, edge betweenness and topic mismatch can flag the edge.

\subsection{Application in Citation and Academic Graphs}
In academic recommendation or paper classification systems, adversaries could inject artificial citations to manipulate topic classification or ranking. STRUC-GUARD+ detects such fake citations using semantic topic mismatch and centrality-aware pruning.

% ==========================================================
\section{Section 11: Future Modifications}
% ==========================================================

\subsection{Heterogeneous Graph Extension}
Future work can extend STRUC-GUARD+ to heterogeneous graphs containing multiple node and edge types. For example, a cybersecurity graph may contain users, devices, IP addresses, files, and processes. The defense can use type-specific topic sets and type-aware message passing:

$$
h_u^{(l+1)}=\sigma\left(\sum_{r\in R}\sum_{v\in N_r(u)}\omega_{uv}^{r}W_rh_v^{(l)}\right)
$$

\subsection{Real-Time Dynamic Temporal Defense}
The Elliptic setting can be extended into a real-time streaming defense. Instead of defending after a snapshot is complete, STRUC-GUARD+ can update suspicious-node scores as transactions arrive:

$$
D_t(u)=\lambda D_{t-1}(u)+(1-\lambda)||x_u^{(t)}-\mu_{t-1}||
$$

This would be useful for online fraud detection.

\subsection{Learned Ontology Rules}
Currently, ontology rules use handcrafted topic sets and thresholds. Future work can learn ontology rules using contrastive learning or graph transformers. A learned semantic consistency function can be:

$$
S(u,v)=MLP([x_u,x_v,|x_u-x_v|,x_u\odot x_v])
$$

\subsection{Quantum-Inspired or Bose-Einstein Graph Learning}
The Bose-Einstein fitness metric can be extended into a quantum-inspired graph regularizer. Nodes can be treated as occupying energy states based on degree and semantic compatibility:

$$
P(u,v)\propto \frac{1}{e^{(E_{uv}-\mu)/T}-1}
$$

This may inspire future graph defense mechanisms that model unnatural edge formation as an energy-based anomaly.

\subsection{Certified Robustness Bounds}
Another future direction is to provide certified guarantees. For example, prove that prediction remains unchanged if the number of edge flips is less than $k$:

$$
f_\Theta(G)=f_\Theta(G+\Delta G)\quad \text{for all } ||\Delta G||_0\leq k
$$

This would make the defense mathematically stronger for high-stakes applications.

\end{document}
